\documentclass[10pt,a4paper]{amsart}
\usepackage[margin=19mm]{geometry}
\usepackage{amsmath,amssymb,mathtools}
\usepackage[scr=rsfs,cal=euler]{mathalfa}
\usepackage{xfrac}
\usepackage[unicode,hidelinks,bookmarks=false]{hyperref}
\newcommand{\ie}{i.e.\ }
\newcommand{\cf}{cf.\ }
\newcommand{\resp}{respectively\ }
\newcommand{\Subsection}{\subsection}
\providecommand{\nobreakdash}{\nobreak\hbox{-}}
\setlength{\emergencystretch}{2em}
\multlinegap0pt
\usepackage{bm}
\usepackage{bbm}
\usepackage{dsfont}
\usepackage{xspace}
\usepackage{cancel}
\usepackage{mathtools}
\usepackage{amsthm}
\makeatletter
\@ifundefined{theorem}{\newtheorem{theorem}{Theorem}}{}
\@ifundefined{lemma}{\newtheorem{lemma}[theorem]{Lemma}}{}
\makeatother
\makeatletter
\newcommand{\bL}{\stackrel{_\leftarrow}{\mathtt L}}
\newcommand{\cH}{\stackrel{_\rightarrow}{\mathscr H}}
\newcommand{\cS}{\ensuremath{\mathscr S}}
\newcommand{\cbH}{\stackrel{_\leftarrow}{\mathscr H}}
\newcommand{\dd}{{\mathrm d}}
\expandafter\ifx\csname ecn\endcsname\relax
\newenvironment{ecn}{\begin{equation}}{\end{equation}}
\fi
\expandafter\ifx\csname esn\endcsname\relax
\newenvironment{esn}{\begin{equation*}}{\end{equation*}}
\fi
\newcommand{\fL}{\stackrel{_\rightarrow}{\mathtt L}}
\newcommand{\leb}{\text{Leb}}
\expandafter\ifx\csname listeo\endcsname\relax
\newenvironment{listeo}{\begin{list}{\alph{enumi})}{\usecounter{enumi}}}{\end{list}}
\fi
\newcommand{\ltwop}[2]{\ensuremath{\!\left\langle{#1,#2}\right\rangle}}
\expandafter\ifx\csname numerai\endcsname\relax
\newenvironment{numerai}{\begin{list}{\roman{enumi})}{\usecounter{enumi}}}{\end{list}}
\fi
\renewcommand{\vec}[1]{\ensuremath{\boldsymbol{#1}}}
\makeatother
\title[A Unified Framework from Boltzmann Transport to Proton Treat]{A Unified Framework from Boltzmann Transport to Proton Treatment Planning}
\author{Andreas E. Kyprianou, Aaron Pim, Tristan Pryer}
\date{Source: arXiv:2508.10596v1 (CC BY 4.0)}
\begin{document}
\maketitle

\noindent\emph{Source and licence.} A Unified Framework from Boltzmann Transport to Proton Treatment Planning. Version arXiv:2508.10596v1 (CC BY 4.0), distributed under the Creative Commons Attribution 4.0 International licence, as recorded in the arXiv licence metadata for this exact version and shown on the abstract page https://arxiv.org/abs/2508.10596v1. Full rights notice: licenses/arxiv-2508.10596-CC-BY-4.0.txt.\par\smallskip

\noindent\textit{Editorial scope.} Counted unit: the Lemma whose source label is \texttt{lemma:monotone} in the article (source0.tex lines 1286--1322, source label \texttt{lemma:monotone}), delivered together with the article's own complete original proof of it (source0.tex lines 1324--1370, one \texttt{proof} environment of 47 lines). No mathematics is written, repaired or re-derived: statement and proof are the article's own text line for line. Required context retained uncounted: eq:pi\_is\_prob\_dist (lines 1044--1049). Every definition, display and label of those lines is retained unchanged. Omitted from this package and not redistributed: 1--1043, 1050--1285, 1323--1323, 1371--2831. Nothing inside a retained excerpt is omitted or altered: every retained body is delivered exactly as the article's lines are written, so no omission receipt of any kind accompanies this package. Citations: the retained excerpts carry no citation command at all, so no bibliography entry of the article is transcribed and no reference is redistributed. Reference adaptation: none was needed and none was made; every reference inside the delivered unit resolves inside it, so no unresolved reference or citation is produced. Printed slips kept literal: no printed word, symbol, hypothesis or number of the counted statement or of its proof was altered. Layout and class: the article's own class is \texttt{amsart} with packages including amsmath, amsthm, todonotes, amsopn, amssymb, eurosym, amsfonts, vmargin, multirow, float, geometry, booktabs, algorithm, algpseudocode; the wrapper is the lane's pinned \texttt{amsart} editorial wrapper at 10pt on A4, which restates the theorem-like environments the retained text needs and adds the article's own macro definitions that the retained text needs (\texttt{\textbackslash bL}, \texttt{\textbackslash cH}, \texttt{\textbackslash cS}, \texttt{\textbackslash cbH}, \texttt{\textbackslash dd}, \texttt{\textbackslash fL}, \texttt{\textbackslash leb}, \texttt{\textbackslash ltwop}, \texttt{\textbackslash vec}), with extra wrapper packages (bm, bbm, dsfont, xspace, cancel, mathtools). The delivered numbering is the unit's own. Assets: no retained span contains a figure environment, a graphics-inclusion command, a PDF-page inclusion, a table or a TikZ picture; no image file is copied into this package, the package declares no asset dependency and no image file is redistributed with it. Licence: version arXiv:2508.10596v1 of this article is distributed under the Creative Commons Attribution 4.0 International licence, as recorded in the arXiv licence metadata for this exact version and shown on the abstract page https://arxiv.org/abs/2508.10596v1. A full rights notice is delivered at licenses/arxiv-2508.10596-CC-BY-4.0.txt. Characters: every retained excerpt is pure ASCII, so the delivered engine, XeLaTeX with its default Latin Modern text font, reports no missing character.
\begin{equation}
  \int_{\mathbb{S}^2} \int_{\mathcal{E}} 
  \pi(\vec{x}, \omega', E';  \omega,  E)\dd \omega\dd E = 1,
  \qquad \text{for all } (\vec{x}, \omega', E') \in \mathcal{C},
  \label{eq:pi_is_prob_dist}
\end{equation}

\begin{lemma}[Monotonicity of the BFP Operator]
  \label{lemma:monotone}
  Assume that the angular diffusion coefficient $\mu \in
  L^{\infty}(D \times \mathcal{E}; \mathbb{R}^+)$, and that the
  stopping power $S \in W^{1,\infty}(D \times \mathcal{E})$
  satisfies
  \begin{equation}
    \lim_{E \downarrow E_{\min}} S(\vec x, E) \geq 0,~ \forall \vec x \in D, 
    \qquad 
    \partial_E S(\vec x, E) \leq 0,~ \forall (\vec x, E) \in D \times \mathcal{E}.
  \end{equation}
  Additionally, suppose that the nuclear scattering cross-section satisfies
  $\sigma_{\mathrm{n}}(\vec{x}, \omega, E) \geq 0$, and that the scattering kernel
  $\pi$ satisfies the normalisation condition \eqref{eq:pi_is_prob_dist}.
  Then, there exist constants $C_{\rightarrow}, C_{\leftarrow} > 0$ such that the
  Boltzmann--Fokker--Planck operator $\fL$ and its adjoint $\bL$ are
  monotone in the sense that
  \begin{equation}
    \begin{aligned}
        \ltwop{\fL u}{u} 
        &\geq
        C_{\rightarrow} \left(\|u\|^2_{\Gamma^+} 
        + \|\sqrt{\mu}\nabla_{\omega}u\|^2_{\leb2(\mathcal{C})} 
        + \|u\|^2_{\leb2(\mathcal{C})} 
        + \lim_{E \downarrow E_{\min}} \|\sqrt{S}u\|^2_{\leb2(D\times \mathbb{S}^2)} \right),
        \quad \forall u \in \cH_0^-, 
    \\ 
        \ltwop{\bL \phi}{\phi} 
        &\geq 
        C_{\leftarrow} \left(\|\phi\|^2_{\Gamma^-} 
        + \|\sqrt{\mu}\nabla_{\omega}\phi\|^2_{\leb2(\mathcal{C})} 
        + \|\phi\|^2_{\leb2(\mathcal{C})} 
        + \lim_{E \uparrow E_{\max}} \|\sqrt{S}\phi\|^2_{\leb2(D\times \mathbb{S}^2)} \right), 
        \quad \forall \phi \in \cbH_0^+.
    \end{aligned}
  \end{equation}
\end{lemma}

\begin{proof}
We first consider the transport term. Integration by parts yields
\begin{equation}
  \ltwop{\omega \cdot \nabla_{\vec x} u}{u}
  =
  \frac{1}{2} \|u\|^2_{\Gamma^+}
  - \frac{1}{2} \|u\|^2_{\Gamma^-}.
\end{equation}
Since $u \in \cH_0^-$ vanishes on $\Gamma^-$, it follows that
\begin{equation}
  \ltwop{\omega \cdot \nabla_{\vec x} u}{u}
  = \frac{1}{2} \|u\|^2_{\Gamma^+} \geq 0.
\end{equation}
Similarly, for the adjoint operator,
\begin{equation}
  -\ltwop{\omega \cdot \nabla_{\vec x} \phi}{\phi}
  = \frac{1}{2} \|\phi\|^2_{\Gamma^-} \geq 0,
  \quad \forall \phi \in \cbH_0^+.
\end{equation}
For the angular diffusion term
\begin{equation}
  -\ltwop{\mu \Delta_\omega v}{v}
  = \|\sqrt{\mu} \nabla_\omega v\|^2_{\leb2(\mathcal{C})} \geq 0,
  \quad \forall v \in \cH_0^- \cup \cbH_0^+.
\end{equation}
We next consider the forward energy term
\begin{align}
  -\ltwop{\partial_E(S u)}{u}
  &= - \ltwop{S \partial_E u}{u} - \ltwop{(\partial_E S) u}{u} \\
  &= \frac{1}{2} \lim_{E \downarrow E_{\min}} \|S^{1/2} u\|^2_{L^2(D \times \mathbb{S}^2)}
  - \frac{1}{2} \ltwop{(\partial_E S) u}{u}.
\end{align}
By the assumption $\partial_E S \leq 0$ and $S \geq 0$ near
$E_{\min}$, both terms are nonnegative. For the backward energy term
\begin{equation}
  \ltwop{S \partial_E \phi}{\phi}
  = - \frac{1}{2} \ltwop{\partial_E S \cdot \phi}{\phi} \geq 0.
\end{equation}
Finally, for the scattering terms
\begin{equation}
  \ltwop{(\sigma_{\mathrm{n}} - \cS) v}{v}
  = \int_{\mathcal{C}} \left( \sigma_{\mathrm{n}} |v|^2 - \cS v \cdot v \right) \geq 0,
\end{equation}
by the positivity and contraction property of $\cS$ (due to
normalisation of $\pi$) and similarly for $\bL$ and $\cS^*$. This
completes the proof.
\end{proof}

\end{document}
